
\section{RIT}
\label{section:rit}
\begin{frame}%STARTCONTENT
\begin{itemize}
  \item Hört sich in SSB die Gegenstation zu hoch oder tief an $\rightarrow$ \emph{Empfangsfrequenz} korrigieren
  \item Dazu den Regler \emph{RIT} (Receiver Incremental Tuning) verwenden
  \item Beim Verstellen der \emph{Sendefrequenz} versteht mich die Gegenstation schlecht
  \item Ändern beide Stationen jedes Mal die Sendefrequenz entsteht \enquote{über das Band wandern}
  \end{itemize}

\end{frame}

\begin{frame}
\only<1>{
\begin{QQuestion}{NF111}{Die Gegenstation antwortet Ihrem Anruf in SSB-Sprechfunk auf einer geringfügig abweichenden Frequenz. Was sollten Sie tun, um die Empfangsfrequenz Ihres Transceivers anzupassen, ohne dabei die Sendefrequenz zu verstellen?}{Passband-Tuning verstellen}
{Notchfilter einschalten}
{RIT-Einstellung ändern}
{PTT betätigen}
\end{QQuestion}

}
\only<2>{
\begin{QQuestion}{NF111}{Die Gegenstation antwortet Ihrem Anruf in SSB-Sprechfunk auf einer geringfügig abweichenden Frequenz. Was sollten Sie tun, um die Empfangsfrequenz Ihres Transceivers anzupassen, ohne dabei die Sendefrequenz zu verstellen?}{Passband-Tuning verstellen}
{Notchfilter einschalten}
{\textbf{\textcolor{DARCgreen}{RIT-Einstellung ändern}}}
{PTT betätigen}
\end{QQuestion}

}
\end{frame}

\begin{frame}
\only<1>{
\begin{QQuestion}{NF112}{Ihr QSO-Partner meldet bei einem SSB-QSO im \qty{2}{\m}-Band: \glqq Sie senden nicht exakt auf meiner Frequenz\grqq{}. Was könnte die Ursache sein?}{CW-Filter ist aktiviert.}
{LSB ist eingestellt.}
{USB ist eingestellt.}
{RIT ist aktiviert.}
\end{QQuestion}

}
\only<2>{
\begin{QQuestion}{NF112}{Ihr QSO-Partner meldet bei einem SSB-QSO im \qty{2}{\m}-Band: \glqq Sie senden nicht exakt auf meiner Frequenz\grqq{}. Was könnte die Ursache sein?}{CW-Filter ist aktiviert.}
{LSB ist eingestellt.}
{USB ist eingestellt.}
{\textbf{\textcolor{DARCgreen}{RIT ist aktiviert.}}}
\end{QQuestion}

}
\end{frame}%ENDCONTENT
